% concatenated from moebius_main_shen.tex
\documentclass[final,12pt]{elsarticle}
\usepackage{xcolor,soul,cancel}
\usepackage[color]{showkeys}
\usepackage{amssymb,amsmath,amsthm,mathtools}
\usepackage{graphicx,float,enumitem,hyperref}
\newlist{alphalist}{enumerate}{1} 
\setlist[alphalist]{label=(\alph*)} 
\journal{arXiv}
\definecolor{refkey}{rgb}{0,1,1}
\definecolor{labelkey}{rgb}{1,0,0}

\numberwithin{equation}{section}

\usepackage[a4paper,margin=27mm]{geometry}
\usepackage{amsmath,amssymb,amsthm,mathtools}
\usepackage{microtype}
\usepackage{enumitem}

\newtheorem{thm}{Theorem}[section]
\newtheorem{lem}[thm]{Lemma}
\newtheorem{cor}[thm]{Corollary}
\newtheorem{prop}[thm]{Proposition}
\newtheorem{rem}[thm]{Remark}
\newtheorem{ex}[thm]{Example}
\newtheorem{defi}[thm]{Definition}
\newtheorem{conj}[thm]{Conjecture}
\newtheorem{theorem}{Theorem}
\newtheorem{lemma}[theorem]{Lemma}
\newtheorem{proposition}[theorem]{Proposition}
\theoremstyle{remark}
\newtheorem{remark}[theorem]{Remark}

\newcommand{\R}{\mathbb R}
\newcommand{\E}{\mathbb E}
\newcommand{\Var}{\operatorname{Var}}
\newcommand{\Cov}{\operatorname{Cov}}
\newcommand{\Unif}{\operatorname{Unif}}
\newcommand{\dd}{\,d}
\newcommand{\Gini}{\mathcal D}

\theoremstyle{definition}
\newtheorem{definition}{Definition}




\begin{document}
\begin{frontmatter}
\title{The Zaporozhets–Tarasov Conjecture on Mean Distances}
\author[eric]{Eric Shen}
\ead{erick.2013@yandex.ru}
\address[eric]{Moscow State University, Moscow, 119991, Russia \\
Institute for Numerical Mathematics, Russian Academy of Sciences, Russia}
\date{August 2026}

\begin{abstract}
For a convex body  $K\subset\R^d$  let $\Delta(K)$ be the expected distance between two independent uniform points of $K$, and let $\theta(K)$ be the expected distance between two uniform points on $\partial K$. The Zaporozhets–Tarasov conjecture asserts $\Delta(K) \le \theta(K)$. We prove this conjecture in case $d=2$. In addition, we give a six-vertex convex polytope in $\mathbb R^3$ for which the reverse strict inequality holds, and obtain counterexamples in every dimension $d \ge 3$ by taking products with segments. This settles the conjecture in every dimension $d \geq 2$. Finally, we show that
$\theta(K) > \frac{\operatorname{per}K}{6}$ for every planar convex body.
\\ \ \\
MSC2020: 52A40, 60D05, 52A22
\end{abstract}

\begin{keyword}
convex body; mean distance; Gini mean difference; Zaporozhets--Tarasov conjecture.
\end{keyword}
\end{frontmatter}

\section{Introduction}
Let $K \subset \mathbb{R}^d$ be a compact convex body. The conjecture of Zaporozhets and Tarasov asks whether
$$
\Delta(K):=\E|X_1-X_2|\leq \E|B_1-B_2|=: \theta(K)
$$
for every convex body $K\subset\R^d$, where $X_1,X_2$ are uniform in $K$ with respect to the normalized Lebesgue measure and $B_1,B_2$ are uniform on $\partial K$ with respect to the normalized arclength measure. 


Lotnikov proved the planar centrally symmetric case and an analogue for sufficiently high moments in arbitrary dimension \cite{Lotnikov}. Tokmachev proved the conjecture, in fact a convex-order statement, for planar bodies whose area and boundary centroids coincide and for a class of circumscribed polytopes \cite{Tokmachev2026}. Related extremal questions for $\Delta$ and $\theta$ were studied in \cite{BonnetEtAl,TokmachevBoundary}.
In particular, Bonnet, Gusakova, Th\"ale and Zaporozhets proved the sharp inequalities
$$
    \frac{7}{60}
    <
    \frac{\Delta(K)}{\operatorname{per} K}
    <
    \frac{1}{6}
$$
for planar convex bodies~\cite{BonnetEtAl}, where $\operatorname{per} K$ stands for perimeter of $K$.  It is also shown that the upper bound is approached by thin rectangles converging to a doubled segment.
According to Tokmachev~\cite[Eq.~(5)]{Tokmachev2026},
Gusakova and Zaporozhets suggested an analogous pair of inequalities for $\theta$
$$
    \frac{1}{6}
    <
    \frac{\theta(K)}{\operatorname{per} K}
    \leq
    \frac{2}{\pi^2}.$$
The upper bound, with equality only for the disk, was proved by
Tokmachev~\cite{TokmachevBoundary}; we prove the lower bound, see Theorem~\ref{thm:lower-boundary-distance}.
\par Prior to that, we prove the Zaporozhets-Tarasov conjecture in the case $d=2$, see Corollary~\ref{corol:conjecture_res}. The result is mostly self-contained, in particular, it does not use the aforementioned inequalities on $\Delta(K)$ from \cite{BonnetEtAl} and is obtained as a corollary of a more general statement, see Theorem~\ref{thm:endpoint}. We also prove an analogous inequality for the second moment, see Corollary~\ref{cor:sec_mom_planar}. In addition to that, we show that the conjecture is false in every dimension $d\geq3$, see Example~\ref{ex:counter-ex}, Figure~\ref{fig:body} and Proposition~\ref{prop:counter-ex}.

\begin{defi}
    For a probability measure $\rho$ on $\R$, write
$$
 \Gini(\rho)=\iint_{\R^2}|x-y|\,d\rho(x)d\rho(y).
$$
This is the \textit{Gini mean difference} of $\rho$.
\end{defi}
We slightly abuse this notation throughout the text, writing $\Gini(X)$ with $X$ being a random variable or a probability density. In either case, $\Gini(X)$ is understood as $\Gini(\rho_X)$, where $\rho_X$ denotes the corresponding probability measure. For basic properties of the Gini mean difference used throughout the paper we refer the reader to \cite{Yitzhaki}.
The main result immediately follows from a stronger directional inequality.

\begin{thm}\label{thm:main}
Let $K\subset\mathbb R^2$ be a convex body. Fix a unit vector $u$. Let $p_u$ be the pushforward of normalized area measure on $K$ under $x\mapsto\langle x,u\rangle$, and let $q_u$ be the corresponding pushforward of normalized arclength measure on $\partial K$. Then
\begin{equation}\label{eq:directional}
 \Gini(q_u)\ge \Gini(p_u).
\end{equation}
\end{thm}

We first prove the theorem for smooth strictly convex bodies and the general case follows by approximation.

\medskip
\noindent\textbf{Plan.}
After rotating and rescaling, the projection of $K$ is $[0,1]$ and
$$
 K=\{(x,y):0\le x\le1,\ g(x)\le y\le f(x)\},
$$
where $f$ is concave and $g$ is convex. Put
$$
 W=f-g,\qquad a=f',\qquad b=g'.
$$
The projected area density is proportional to $W$, while the projected boundary density is proportional to
$$
 \sqrt{1+a^2}+\sqrt{1+b^2}.
$$
The proof has two steps. First, we find two densities $\varphi$ and $\psi$, nonincreasing and nondecreasing, respectively, such that the projected arclength density $q$ satisfies $q = \frac{\varphi+\psi}{2}$ while the projected area density $p$ is proportional to the difference of the CDFs of $\varphi$ and $\psi$. Then we apply a one-dimensional Theorem saying that an endpoint mixture is more dispersed than the normalized occupation measure between the endpoints.
The only substantial step is the second one, to which we now turn.

\section{The monotone-endpoint theorem}

\begin{thm}
    \label{thm:endpoint}
Let $\varphi$ be a nonincreasing probability density and $\psi$ a nondecreasing probability density on $[0,1]$. Let
$$
 \Phi(x)=\int_0^x\varphi(s)\,ds,
 \qquad
 \Psi(x)=\int_0^x\psi(s)\,ds.
$$
Put
$$
 H=\Phi-\Psi,
 \qquad
 A=\int_0^1H(x)\,dx,
$$
and assume $A>0$. Define
$$
 p(x)=\frac{H(x)}A,
 \qquad
 q(x)=\frac{\varphi(x)+\psi(x)}2.
$$
Then
$$
 \Gini(q)\ge \Gini(p).
$$
\end{thm}

 Since $\Phi$ is concave and $\Psi$ is convex, $\Phi(x)\ge x \ge \Psi(x)$.
 Thus $H\ge0$, and $p$ is indeed a probability density.

\begin{rem}
There is also a geometric reformulation.  Let $x(u)=\Phi^{-1}(u)$ and $y(u)=\Psi^{-1}(u)$. If $U\sim\Unif[0,1]$, then $ x(U)\sim \varphi,\,\,  y(U) \sim \psi$. Consider
$$
    L=\{(u,z):0\leq u\leq 1,\ x(u)\leq z\leq y(u)\}.
$$
Since $x$ is convex and $y$ is concave, $L$ is convex. Choosing one of the two endpoints
of the (uniformly distributed) vertical slice $[x(U),y(U)]$ with equal probability gives the law $q$.
Schematically,
$$
 q(z)dz=\frac12\int_0^1\bigl(\delta_{x(u)}+\delta_{y(u)}\bigr)(dz)\,du.
$$
 On the other hand, the length of the horizontal section of $L$ at height $z$ is
    $$\bigl|\{u:x(u)\leq z\leq y(u)\}\bigr|
    =\Phi(z)-\Psi(z)=H(z). $$
Consequently, $p$ is the second-coordinate distribution of a uniform
point of $L$, that is, choose a slice with a probability proportional to its length and then a random point in it. Explicitly, 
$$
 p(z)\,dz
 =\frac{\displaystyle\int_0^1
 \mathbf 1_{\{x(u)\le z\le y(u)\}}\,du}
 {\displaystyle\int_0^1(y(u)-x(u))\,du}\,dz.
$$ 
\end{rem}

Every decreasing probability density on $[0,1]$ is a mixture of uniform densities on intervals $[0,t]$, as shown by the following Proposition. This is a one-sided form of Khintchine’s representation theorem for unimodal distributions, see \cite{Khintchine}.

\begin{prop} \label{prop:monotone_factoriz}
Let $\varphi$ be a probability density on $[0,1]$ admitting a nonincreasing representative, which we also denote by $\varphi$. There is a random variable $T\in[0,1]$ such that, for $U\sim\Unif[0,1]$ independent of $T$, the variable
$$
 X=TU
$$
has density $\varphi$. Likewise, for every nondecreasing density $\psi$, there is $S\in[0,1]$ such that
$$
 Y=1-SV,
 \qquad V\sim\Unif[0,1],
$$
has density $\psi$.
\end{prop}
    \begin{proof}
        
     As the law of $T$ one may take the Stieltjes mixing measure 
$$
 \rho(dt)=-t\,d\varphi(t)+\varphi(1-)\delta_1(dt),
$$
for which
$$
 \varphi(x)=\int_{[x,1]}\frac1t\,\rho(dt) \quad \text{for a.e. } x\in (0,1) .
$$ 
As $\varphi$ is nonincreasing, $\rho$ is nonnegative and the Stieltjes integration by parts gives $\rho ([0,1]) = \int_0^1 \varphi = 1.$
\end{proof}

Next lemma gives a lower bound on the Gini mean difference of $q$ in terms of its mean and auxiliary variables.

\begin{lem}\label{lem:low_bound_enpdpoint}
    In the notation of Theorem~\ref{thm:endpoint} and Proposition~\ref{prop:monotone_factoriz}, denote 
    $$ A=\int_0^1(\Phi-\Psi)
   =\E Y-\E X
   =1-\frac{\E T+\E S}{2}.$$
   Then $\Gini(q) \geq \frac{1}{3} B := \frac{1}{3}(1-A+\frac{3}{2}A^2 + \Var(T) + \Var(S)).$
\end{lem}
\begin{proof}
    

Take $U, U'\sim \Unif[0,1]$ independent. Then for $s,t\in[0,1]$, direct integration gives

$$ \E|tU-sU'|
 =\frac{t+s}{6}
 +\frac{(t-s)^2}{3\max(t,s)},$$

with the second term read as zero when $s=t=0$. Since $\max(s,t)\le1$,
$$
 \E|tU-sU'|
 \ge \frac{t+s}{6}+\frac{(t-s)^2}{3}.
$$
Now  invoke Proposition~\ref{prop:monotone_factoriz}:
$$\Gini(X) = \E|TU-T'U'| = \E_{T,T'}\E(|TU-T'U'| \, |\,T, T') \geq \E(\frac{T+T'}{6}+\frac{(T-T')^2}{3}).$$
Using that $T$ and $T'$ are independent and applying the same argument to $Y$ one obtains
\begin{equation}\label{eq:Dmunu}
 \Gini(X)\ge\frac{\E T}{3}+\frac23\Var T,
 \qquad
 \Gini(Y)\ge\frac{\E S}{3}+\frac23\Var S.
\end{equation}

For independent laws on the line with finite first moments with distribution functions $\Phi,\Psi$,
$$ 2\E|X-Y|-\Gini(X)-\Gini(Y)
 =2\int_{\R}(\Phi-\Psi)^2.$$

This follows immediately from the standard CDF formulas for the three terms. Since $q$ is an equal-probability mixture of $X$ and $Y$, we have 
\begin{equation}\label{eq:Dqidentity}
 \Gini(q)
 =\frac{\Gini(X)+\Gini(Y)}2
  +\frac12\int_0^1H(x)^2\,dx.
\end{equation}
By Cauchy--Schwarz, $\int H^2\ge A^2$. Finally, using \eqref{eq:Dmunu}, and \eqref{eq:Dqidentity}, we obtain
\begin{equation}\label{eq:Dqlower}
 \Gini(q)\ge\frac{1}{6} (\E T + \E S + 2 \Var T + 2 \Var S) + \frac{1}{2}A^2 = \frac{1}{3}B.
\end{equation}
\end{proof}


\begin{lem} \label{lem:upper_bound} Let $p$ be as in Theorem~\ref{thm:endpoint}, let $Z$ have density $p$, and assume $B$ as in Lemma~\ref{lem:low_bound_enpdpoint}. Then
$$ 12\Var Z\le B^2.$$

\end{lem}
If $P$ denotes the CDF of $Z$, then  by the Gini covariance identity

$$ \Gini(p)=4\Cov(Z,P(Z)).$$
Since $p$ is absolutely continuous, $P(Z)$ is uniform on $[0,1]$. Consequently, using Cauchy--Schwarz again,

$$ \Gini(p) \le 4 \sqrt{\Var(Z) \Var(P(Z)} = 4\sqrt{\Var(Z) \frac{1}{12}}
 = \frac2{\sqrt3}\sqrt{\Var Z}.$$

Hence with \eqref{eq:Dqlower}, Lemma~\ref{lem:upper_bound} finishes the proof of Theorem~\ref{thm:endpoint}.

\begin{proof}[Proof of Lemma~\ref{lem:upper_bound}]
This part is rather technical, though quite elementary. For every integer $m\ge0$, Fubini's theorem gives
$$
 \E Z^m
 =\frac{\E Y^{m+1}-\E X^{m+1}}{(m+1)A}.
$$
Indeed, 
$$
 \int_0^1x^m\Phi(x)\,dx = \int_0^1x^m\E\mathbf 1_{\{X\leq x\}}\,dx
 =\E\int_X^1x^m\,dx
 =\frac{1-\E X^{m+1}}{m+1},
$$


and since $p=(\Phi-\Psi)/A$, the analogous identity for $\Psi$ gives the claim after subtraction. What follows is mostly algebraic identities involving $T$ and $S$.

 Recall that $X = TU$, and $Y = 1 -SV$, and introduce
$$
 z=\frac{\E T-\E S}{2},\quad
 e=\E[T(1-T)]+\E[S(1-S)],\quad
 k=\E[T(1-T)]-\E[S(1-S)].
$$


Using the general moment formula with $m=1$, and using
$$
 \E X^2=\frac{\E[T^2]}{3},
 \qquad
 \E Y^2
 =1-\E S+\frac{\E[S^2]}{3},
$$
we obtain
$$
 \E Z
 =
 \frac{
  1-\E S+\frac13\bigl(\E[S^2]-\E[T^2]\bigr)
 }{2A} =\frac12+\frac{z+k}{6A}.
$$

Similarly, using the moment formula with $m=2$,
$$
 \E X^3=\frac{\E[T^3]}4,
$$
while
$$
 \E Y^3
 =
 \E(1-SV)^3
 =
 1-\frac32\E S+\E[S^2]-\frac14\E[S^3].
$$
Hence
$$
 \E Z^2
 =
 \frac{
  1-\frac32\E S+\E[S^2]
  -\frac14\bigl(\E[S^3]+\E[T^3]\bigr)
 }{3A}.
$$
The numerator can be regrouped as
$$
 1-\frac32\E S+\E[S^2]
  -\frac14\bigl(\E[S^3]+\E[T^3]\bigr) =
 A+\frac{z+k}{2}
 -\frac14\left(
   \E[T(1-T)^2]+\E[S(1-S)^2]
 \right). $$
Consequently,
$$
 \E Z^2
 =
 \frac13+\frac{z+k}{6A}
 -\frac{
   \E[T(1-T)^2]+\E[S(1-S)^2]
 }{12A}.
$$

Subtracting $(\E Z)^2$ gives
$$ 12\Var Z =
 12\E Z^2-12(\E Z)^2 = 1 -\frac{\E[T(1-T)^2]+\E[S(1-S)^2]}{A} -\frac{(z+k)^2}{3A^2}.$$

We now estimate the two expectations which occur with a minus sign.
By Cauchy--Schwarz,
$$
 \bigl(\E[T(1-T)]\bigr)^2
 \le
 \E T\;\E[T(1-T)^2].
$$
Since
$$
 \E[T(1-T)]=\frac{e+k}{2},
 \qquad
 \E T=1-A+z,
$$
this gives
$$
 \E[T(1-T)^2]
 \ge
 \frac{(e+k)^2}{4(1-A+z)}.
$$
Likewise,
$$
 \E[S(1-S)^2]
 \ge
 \frac{(e-k)^2}{4(1-A-z)}.
$$
Substitution into the preceding formula yields
$$
 12\Var Z\le 1-\Lambda,
$$
where
$$
 \Lambda=
 \frac{(e+k)^2}{4A(1-A+z)}
 +\frac{(e-k)^2}{4A(1-A-z)}
 +\frac{(z+k)^2}{3A^2}.
$$


We claim that
$$
 \Lambda\ge2(1-B),
$$
where $B=1-A+\frac{3}{2}A^2 + \Var(T) + \Var(S) = 1+A-\frac{1}{2}A^2-2z^2-e$.
\par Once this is proved, the result follows immediately. Indeed, $B\ge0$ from its original definition, and therefore
$$
 12\Var Z\le1-\Lambda\le2B-1\le B^2,
$$

It remains to prove the claim. Since
$\E T=1-A+z$ and $\E S=1-A-z$ lie in $[0,1]$,
$$
 0\le z^2\le\min\{A^2,(1-A)^2\}.
$$
Using the latter formula for $B$, we obtain

$$ \Lambda-2(1-B)
 ={}\frac{(e+k)^2}{4A(1-A+z)}
 +\frac{(e-k)^2}{4A(1-A-z)}
 +\frac{(z+k)^2}{3A^2}
 +2A-A^2-4z^2-2e.$$

This is a convex quadratic polynomial in $k$ and $e$, which we estimate from below in a standard way. Recall $1-A+z = \E T$, $1-A-z=\E S$ and, to simplify the analysis, introduce new variables $e=x+y$, $k=x-y$. Then
$$\Lambda-2(1-B) = F(x,y) = \frac{x^2}{A \E T} + \frac{y^2}{A \E S} + \frac{(z+x-y)^2}{3A^2} + 2A-A^2-4z^2-2x-2y.$$
Minimizing that in $x,y$ gives
$$F(x,y) \geq \frac{A^3(A+2)-(4A-1)z^2}{A(A+2)}.$$

% Derivation of the lower bound for F.
%
% Set $$d=z+x-y.$$
%
% At the critical point,
% $$\frac{x}{A\,\E T}+\frac{d}{3A^2}=1,\qquad\frac{y}{A\,\E S}-\frac{d}{3A^2}=1.$$
%
% Hence
% $$x=A\,\E T-\frac{\E T}{3A}d,\qquad  y=A\,\E S+\frac{\E S}{3A}d.$$
%
% Since
% $$\E T-\E S=2z,\qquad\E T+\E S=2(1-A),$$
% we obtain
% $$d=z+x-y=(1+2A)z-\frac{2(1-A)}{3A}d.$$

% Therefore
% $$d=\frac{3A(2A+1)}{A+2}\,z.$$

% To evaluate $F$ at the critical point without substituting
% $x$ and $y$ directly, multiply the two critical-point equations
% by $x$ and $y$, respectively, and add. This gives
% $$\frac{x^2}{A\,\E T}+\frac{y^2}{A\,\E S}=x+y-\frac{d(d-z)}{3A^2}.$$

% Consequently,
% $$F_{\min} =-(x+y)+\frac{zd}{3A^2} +2A-A^2-4z^2.$$
% $$ x+y  = A(\E T+\E S)  +\frac{\E S-\E T}{3A}d  = 2A(1-A)-\frac{2zd}{3A}. $$
% Thus  $$ F_{\min}  = A^2-4z^2+\frac{(2A+1)zd}{3A^2}.$$
% Substituting the expression for $d$ yields $$
% F_{\min} =A^2-4z^2 +\frac{(2A+1)^2}{A(A+2)}z^2 =\frac{A^3(A+2)-(4A-1)z^2}{A(A+2)}.$$



Thus it only remains to verify that the numerator is nonnegative.






If $0<A\le\frac14$, then $4A-1\le0$, so $$A^3(A+2)-(4A-1)z^2\ge A^3(A+2) \geq 0.$$ If
$\frac14\le A\le\frac12$, then $z^2\le A^2$ and
$$
 A^3(A+2)-(4A-1)z^2\ge A^3(A+2)-(4A-1)A^2
 =A^2(1-A)^2\ge0.
$$
Finally, if $\frac12\le A<1$, then $z^2\le(1-A)^2$ and
$$
 \begin{aligned}
 A^3(A+2)-(4A-1)z^2&\ge A^3(A+2)-(4A-1)(1-A)^2\\
  &=(A^2-A)^2+(4A-1)(2A-1)\ge0.
 \end{aligned}
$$
Therefore $\Lambda-2(1-B)\ge0$, and hence $12\Var Z\le B^2$.

\end{proof}
The analogous result for the second moment is immediate.
\begin{cor}\label{cor:sec_mom}
    In the setting of Theorem~\ref{thm:endpoint}, let $Q$ be a random variable with density $q$, and let $Z$, as before, have density $p$. Then  $\Var(Q) \geq \Var(Z)$.
\end{cor}
\begin{proof}
    Indeed, we have $\frac{1}{3}B \leq \Gini(q) \leq \frac{2}{\sqrt{3}} \sqrt{\Var(Q)}$, therefore $12 \Var(Q) \geq B^2$. At the same time $12 \Var(Z) \leq B^2$, by Lemma~\ref{lem:upper_bound}.
\end{proof}

\section{Proof of Theorem~\ref{thm:main}}

Assume for now that $K$ is smooth and strictly convex. Fix a direction, normalize its projection to $[0,1]$, and use the notation from the plan:
$$
 W=f-g,\qquad a=f',\qquad b=g'.
$$
The functions $a$ and $b$ are respectively nonincreasing and nondecreasing, while
$$
 W'=a-b,\qquad W(0)=W(1)=0.
$$
The normalized projected area density is
\begin{equation}\label{eq:area-density}
 p(x)=\frac{W(x)}{\mathcal A},
 \qquad
 \mathcal A=\int_0^1W(x)\,dx.
\end{equation}
The unnormalized projected arclength boundary density is
$$
 r(x)=\sqrt{1+a(x)^2}+\sqrt{1+b(x)^2}.
$$

Let
$$
R=\int_0^1r(x)\,dx,
$$
so that the normalized projected boundary density is
$
q(x)=\frac{r(x)}{R}.
$
Define
$$
\varphi(x)
=
\frac{
\sqrt{1+a(x)^2}+a(x)
+
\sqrt{1+b(x)^2}-b(x)
}{R}
$$
and
$$
\psi(x)
=
\frac{
\sqrt{1+a(x)^2}-a(x)
+
\sqrt{1+b(x)^2}+b(x)
}{R}.
$$
Clearly,
$$
q=\frac{\varphi+\psi}{2}.
$$
 Since $a$ is nonincreasing and $b$ is nondecreasing, it follows that
$\varphi$ is nonincreasing and $\psi$ is nondecreasing. Both are nonnegative. Moreover,
$$
\int_0^1(a-b)\,dx=W(1)-W(0)=0,
$$
and hence
$$
\int_0^1\varphi(x)\,dx
=
\frac{1}{R}\int_0^1\bigl(r(x)+a(x)-b(x)\bigr)\,dx
=1,
$$
and similarly
$$
\int_0^1\psi(x)\,dx=1.
$$
Thus $\varphi$ and $\psi$ are probability densities of the type appearing in
Theorem~\ref{thm:endpoint}.


Let $\Phi$ and $\Psi$ be their distribution functions. Then
$$
\Phi(x)-\Psi(x)
=
\frac{2}{R}\int_0^x(a(s)-b(s))\,ds
=
\frac{2W(x)}{R}.
$$
Consequently,
$$
\int_0^1\bigl(\Phi(x)-\Psi(x)\bigr)\,dx
=
\frac{2 \mathcal A}{R},
$$
and therefore
$$
\frac{\Phi(x)-\Psi(x)}
{\displaystyle\int_0^1(\Phi-\Psi)}
=
\frac{W(x)}{\mathcal A}
=
p(x).
$$
Since $q=(\varphi+\psi)/2$, Theorem~\ref{thm:endpoint} gives
$$
\Gini(q)\ge \Gini(p).
$$

% For $t\in\R$, define
% $$
%  r_t(x)=|a(x)-t|+|b(x)-t|,
%  \qquad
%  R_t=\int_0^1r_t(x)\,dx,
% $$
% and let $q_t$ be the probability density $r_t/R_t$. Put $d=a-b$ and define

% $$ \varphi_t=\frac{r_t+d}{R_t},
%  \qquad
%  \psi_t=\frac{r_t-d}{R_t}.$$

% The identities

% $$ \frac{r_t+d}{2}=(a-t)_+ +(t-b)_+,
%  \qquad
%  \frac{r_t-d}{2}=(t-a)_+ +(b-t)_+$$

% show that $\varphi_t$ is decreasing, $\psi_t$ is increasing, and both are nonnegative. Moreover, since $\int_0^1d=W(1)-W(0)=0$, both are probability densities.

% Let $\Phi_t,\Psi_t$ be their CDFs. Then
% $$
%  \Phi_t(x)-\Psi_t(x)
%  =\frac2{R_t}\int_0^x(a-b)
%  =\frac{2W(x)}{R_t}.
% $$
%  Therefore, taking into account that $q_t = \frac{\phi_t + \psi_t}{2}$, Theorem~\ref{thm:endpoint} applied to $q_t$ and $p$ implies
% \begin{equation}\label{eq:qt}
%  \Gini(q_t)\ge\Gini(p)
%  \qquad\text{for every }t\in\R.
% \end{equation}

% It remains to return from those absolute-slope densities to Euclidean arclength. For every $s\in\R$,
% \begin{equation}\label{eq:length-mixture}
%  \sqrt{1+s^2}
%  =\frac12\int_{\R}|s-t|\frac{dt}{(1+t^2)^{3/2}}.
% \end{equation}
% To verify this, differentiate twice in $s$: both sides have second derivative $(1+s^2)^{-3/2}$, and their values and first derivatives agree at $s=0$.

% The unnormalized projected Euclidean boundary density is
% $$
%  r(x)=\sqrt{1+a(x)^2}+\sqrt{1+b(x)^2}.
% $$
% By \eqref{eq:length-mixture},
% $$
%  r(x)=\frac12\int_{\R}r_t(x)\frac{dt}{(1+t^2)^{3/2}}.
% $$
% Let $P=\int_0^1r(x)\,dx$ and define
% $$
%  \alpha(dt)=\frac{R_t}{2P(1+t^2)^{3/2}}\,dt.
% $$
% Taking total masses in the preceding identity shows that $\alpha$ is a probability measure. If $q=r/P$ is the projected boundary density, then

% $$ q=\int_{\R}q_t\,\alpha(dt).$$

% The same holds for the CDFs:
% $$ Q = \int_0^xq(s)ds, \quad Q_t = \int_0^x q_t(s)ds,$$
% and by Tonelli 

% \begin{equation}\label{eq:qmixture}
% Q = \int_0^x \int_{\mathbb R} q_t(s) \alpha(dt) ds = \int_{\mathbb R} Q_t \alpha(dt).
% \end{equation}
% For a probability measure $\rho$ on $[0,1]$ with CDF $Q_\rho$,
% \begin{equation}\label{eq:cdf-gini}
%  \Gini(\rho)=2\int_0^1Q_\rho(x)(1-Q_\rho(x))\,dx.
% \end{equation}
% Since $z\mapsto z(1-z)$ is concave and CDFs depend linearly on the measure, Jensen inequality applied to \eqref{eq:qmixture} implies
% $$
%  \Gini(q)\ge\int_{\R}\Gini(q_t)\,\alpha(dt).
% $$
% Using \eqref{eq:qt},
% $$
%  \Gini(q) \ge \Gini(p).
% $$
This proves the directional inequality \eqref{eq:directional} for every smooth strictly convex planar body.

\par Let $K_n$ be smooth strictly convex bodies converging to $K$ in the
Hausdorff metric. Then the normalized area measures on $K_n$ converge
weakly to normalized area measure on $K$, see e.g.
\cite{Schneider}. Likewise, the normalized arclength measures on
$\partial K_n$ converge weakly to normalized arclength measure on
$\partial K$; this follows from the weak continuity of curvature
measures under Hausdorff convergence, see e.g.
\cite{HugSchneider,Schneider}. Consequently, their pushforwards under
$x\mapsto\langle x,u\rangle$ converge weakly as well. Since all these
measures are supported in a common compact interval, the Gini
functionals converge, and passing to the limit proves the result for
$K$.


\begin{rem}
One can avoid the approximation by keeping the vertical boundary faces as atoms at the endpoints of the projection interval. The limiting argument is used here only to keep the main exposition lighter.
\end{rem}
\begin{cor} \label{corol:conjecture_res}
For every planar convex body $K$,
$$
 \E |X_1-X_2|\le \E |B_1-B_2|,
$$
where $X_1,X_2$ are uniform in $K$ and $B_1,B_2$ are uniform on $\partial K$. 
\end{cor}

\begin{proof}
Indeed, for normalized angular measure $\bar\sigma$ on $S^1$,
\begin{equation}\label{eq:cosine}
 |z|=\frac\pi2\int_{S^1}|\langle z,u\rangle|\,d\bar\sigma(u).
\end{equation}
See e.g. \cite{Schneider}. Integrating \eqref{eq:cosine} against the two-point area and boundary distributions, and then using \eqref{eq:directional}, gives the result.
\end{proof}

\begin{cor}\label{cor:sec_mom_planar}
For every planar convex body $K$,
$$
\mathbb E|X_1-X_2|^2\leq \mathbb E|B_1-B_2|^2,
$$
where $X_1,X_2$ are uniform in $K$ and $B_1,B_2$ are uniform on $\partial K$.
\end{cor}

\begin{proof}
For a smooth strictly convex body, the proof of Theorem~\ref{thm:main}
shows that, for every $u\in S^1$, the projected area and boundary densities are
obtained from the same pair of decreasing and increasing densities as in
Theorem~\ref{thm:endpoint}. Corollary~\ref{cor:sec_mom}
gives
$$
\operatorname{Var}\langle X_1,u\rangle
\leq
\operatorname{Var}\langle B_1,u\rangle.
$$
Observe that
$$
\mathbb E\langle X_1-X_2,u\rangle^2
=
2\operatorname{Var}\langle X_1,u\rangle,
$$
and similarly
$$
\mathbb E\langle B_1-B_2,u\rangle^2
=
2\operatorname{Var}\langle B_1,u\rangle.
$$
Therefore
$$
\mathbb E\langle X_1-X_2,u\rangle^2
\leq
\mathbb E\langle B_1-B_2,u\rangle^2
$$
for every $u\in S^1$.

For normalized angular measure $\bar\sigma$ on $S^1$,
$$
|z|^2
=
2\int_{S^1}\langle z,u\rangle^2\,d\bar\sigma(u).
$$
Integrating the preceding directional inequality over $u$ therefore gives
$$
\mathbb E|X_1-X_2|^2
\leq
\mathbb E|B_1-B_2|^2.
$$

The general case follows by the same smooth approximation argument as in the proof
of Theorem~\ref{thm:main}, since variance is continuous under weak convergence
for measures supported in a common compact interval.
\end{proof}



\section{A long-body observation}
In this section, we construct a three-dimensional counterexample  to the Zaporozhets-Tarasov conjecture, see Example~\ref{ex:counter-ex}. This example is then used to construct counterexamples in higher dimensions, see Proposition~\ref{prop:counter-ex}. The
main observation, discussed in Remark~\ref{rmk:why}, is that in dimension
three the area and perimeter profiles of sections orthogonal to a fixed
direction can be controlled almost independently. We exploit this
freedom to make the projected Gini mean difference strictly larger for
the volume measure than for the surface measure. A sufficiently large
stretch in the projection direction then yields the required
counterexample. 

\begin{lem}\label{lem:axis}
Suppose a random point has the form $Z=(LT,U)\in\R\times\R^m$, where $T\in[0,1]$ and the support of $U$ has diameter at most $D$. If $Z_1,Z_2$ are independent copies, then
$$
L\E|T_1-T_2|\leq \E|Z_1-Z_2|\leq L\E|T_1-T_2|+D.
$$
\end{lem}

\begin{proof}
The lower bound follows by projection onto the first coordinate. For the upper bound, use
$$
|(L(T_1-T_2),U_1-U_2)|\leq L|T_1-T_2|+|U_1-U_2|.
$$
\end{proof}

Thus, after division by $L$, mean distances in a long body converge to mean absolute differences of the corresponding axial coordinates.

\begin{ex}\label{ex:counter-ex}
    For $L>0$ and $\varepsilon>0$, define
$$
K_{L,\varepsilon}:=\left\{(Lt,x,y):0\leq t\leq1,\quad |x|\leq t,\quad |y|\leq\varepsilon\left(1-\frac{2t}{3}\right)\right\}.
$$
This is the convex hull of $(0,0,\pm\varepsilon)$ and the four points $(L,\pm1,\pm\varepsilon/3)$.
\end{ex}


\begin{figure}[t]
\centering
\includegraphics[width=0.55\textwidth]{counterexample_body.png}
\caption{The polytope $K_{L,\varepsilon}$ in the rescaled coordinates $(t,x,y/\varepsilon)$. One transverse width grows with $t$, while the other decreases.}
\label{fig:body}
\end{figure}

\begin{thm}\label{thm:explicit}
For $L=1000$ and $\varepsilon=1/100$, one has
$$
\Delta(K_{L,\varepsilon})>\theta(K_{L,\varepsilon}).
$$
\end{thm}

\begin{proof}

The area of the section at level $t$ equals
$
4\varepsilon t\left(1-\frac{2t}{3}\right).
$
Hence the axial coordinate of a uniform interior point has density
$$
p_I(t)=\frac{18}{5}t\left(1-\frac{2t}{3}\right)
$$
and distribution function
$$
F_I(t)=\frac{9t^2-4t^3}{5}.
$$
Using $\Gini(p)=2\int_0^1F(t)(1-F(t))\,dt$, we obtain
$$
\Gini(p_I)=\frac{246}{875}=0.281142857\ldots.
$$
Lemma~\ref{lem:axis} therefore gives
$$
\frac{\Delta(K_{L,\varepsilon})}{L}\geq\frac{246}{875}.
$$

We next compute the axial distribution of normalized surface measure. The two faces $x=\pm t$ have combined area density
$$
4\varepsilon\left(1-\frac{2t}{3}\right)\sqrt{L^2+1},
$$
while the two faces $y=\pm\varepsilon(1-2t/3)$ have combined area density
$$
4t\sqrt{L^2+\frac{4\varepsilon^2}{9}}.
$$
The end face $t=1$ has area $4\varepsilon/3$. Put
$$
a=\varepsilon\sqrt{L^2+1},\qquad
b=\sqrt{L^2+\frac{4\varepsilon^2}{9}}-\frac{2\varepsilon}{3}\sqrt{L^2+1},
$$
$$
D=a+\frac b2+\frac{\varepsilon}{3},\qquad
\alpha=\frac aD,\qquad \beta=\frac{b}{2D}.
$$
For $0\leq t<1$, the boundary axial distribution function is
$$
F_{\partial}(t)=\alpha t+\beta t^2;
$$
the remaining mass is an atom at $t=1$. Consequently its Gini mean difference is exactly
$$
g_{L,\varepsilon}
=\alpha+\frac{2\beta}{3}-\frac{2\alpha^2}{3}-\alpha\beta-\frac{2\beta^2}{5}.
$$
The transverse diameter of $K_{L,\varepsilon}$ is at most $2\sqrt{1+\varepsilon^2}$, so Lemma~\ref{lem:axis} gives
$$
\frac{\theta(K_{L,\varepsilon})}{L}
\leq g_{L,\varepsilon}+\frac{2\sqrt{1+\varepsilon^2}}{L}.
$$
For the values suggested in the statement,
$
g_{1000,1/100}=0.2692731635\ldots
$
and hence
$$
\frac{\theta(K_{1000,1/100})}{1000}
\leq0.2712732636\ldots
<0.2811428571\ldots
\leq\frac{\Delta(K_{1000,1/100})}{1000}.
$$
\end{proof}




The family also has a clean asymptotic criterion. As $L\to\infty$, the boundary axial density tends to the normalization of
$$
t+\varepsilon\left(1-\frac{2t}{3}\right),
$$
and its Gini mean difference is
$$
g_{\partial}(\varepsilon)
=\frac{4(19\varepsilon^2+33\varepsilon+9)}{15(4\varepsilon+3)^2}.
$$
Moreover,
$$
g_{\partial}(\varepsilon)-\frac{246}{875}
=\frac{2(746\varepsilon^2+2694\varepsilon-171)}{2625(4\varepsilon+3)^2}.
$$
Thus every
$$
0<\varepsilon<\varepsilon_*:=\frac{-1347+45\sqrt{959}}{746}=0.062396288\ldots
$$
produces counterexamples for all sufficiently large $L$.

\begin{rem}\label{rmk:why}
    Consider a long body with rectangular transverse sections
$$
Q_t=\prod_{i=1}^{d-1}[-a_i(t),a_i(t)].
$$
Up to constant factors, the axial density of volume measure is
$$
\prod_{i=1}^{d-1}a_i(t),
$$
whereas the lateral surface density is
$$
\sum_{i=1}^{d-1}\prod_{j\neq i}a_j(t).
$$
In dimension three this is the difference between a product $a_1a_2$ and a sum $a_1+a_2$. Choosing one width increasing and one decreasing can make the normalized surface density more concentrated along the long axis than the normalized volume density.  In dimension two there is only one transverse width: the second expression is the empty product. 
\end{rem}



\begin{prop}\label{prop:counter-ex}
There are counterexamples in every dimension $d\geq3$.
\end{prop}
\begin{proof}
The case $d=3$ was established above. Let $d=3+m$, where $m\ge 1$, and put
$$
K_L:=K_{L,1/100},
\qquad
C_M:=[-M,M]^m,
\qquad
\widetilde K_{L,M}:=K_L\times C_M .
$$

Since $1/100<\varepsilon_*$, the preceding asymptotic computation shows that there exists
$c>0$ such that, for all sufficiently large $L$,
$$
\Delta(K_L)-\theta(K_L)\ge cL.
$$

The boundary of $\widetilde K_{L,M}$, up to a set of surface measure zero, is the disjoint union
$$
(\partial K_L\times C_M)\cup(K_L\times\partial C_M).
$$
The first part has surface measure
$$
\mathcal H^2(\partial K_L)(2M)^m,
$$
whereas the second has surface measure
$$
|K_L|\,\mathcal H^{m-1}(\partial C_M)
=
2m\,|K_L|(2M)^{m-1}.
$$
Since both $|K_L|$ and $\mathcal H^2(\partial K_L)$ are of order $L$, the relative surface
mass of $K_L\times\partial C_M$ is
$$
\eta_{L,M}=O(M^{-1}),
$$
uniformly for sufficiently large $L$.

Let $\mu_{L,M}$ be normalized surface measure on
$\partial K_L\times C_M$. It is the product of normalized surface measure on
$\partial K_L$ and normalized volume measure on $C_M$. If $Z_1,Z_2$ are independent with law $\mu_{L,M}$, then
$$
\mathbb E|Z_1-Z_2|
\le \theta(K_L)+2M\sqrt m.
$$

The diameter of $\widetilde K_{L,M}$ is $O(L+M)$. Replacing normalized surface
measure on $\partial\widetilde K_{L,M}$ by $\mu_{L,M}$ changes the corresponding mean
distance by at most
$$
O\!\left(\eta_{L,M}(L+M)\right)
=
O\!\left(\frac{L}{M}+1\right).
$$
Consequently,
$$
\theta(\widetilde K_{L,M})
\le
\theta(K_L)
+
O(M)
+
O\!\left(\frac{L}{M}+1\right).
$$
Choose $M=M(L)$ so that
$$
M\longrightarrow\infty,
\qquad
\frac{M}{L}\longrightarrow0
\qquad (L\to\infty),
$$
then the error terms above are $o(L)$.
On the other hand, projection onto the first three coordinates does not increase distances, so
$$
\Delta(\widetilde K_{L,M})\ge \Delta(K_L).
$$

 Combining all inequalities above, we obtain
$$
\Delta(\widetilde K_{L,M})-\theta(\widetilde K_{L,M})
\ge
cL-o(L)>0
$$
for all sufficiently large $L$. Thus $\widetilde K_{L,M}$ is a counterexample in
dimension $d$.
\end{proof}
\section{The lower bound for boundary mean distance}

We now prove the lower estimate stated in the Introduction.  The argument is again a
comparison of one-dimensional distributions.  Here the relevant laws are obtained by
choosing either one random tangent direction or the median of three random tangent
directions on a moving half-boundary.

\begin{thm}\label{thm:lower-boundary-distance}
For every planar convex body $K$,
$$
    \theta(K)>\frac{\operatorname{per}K}{6}.
$$
The constant $1/6$ is sharp.
\end{thm}

Assume first that $\partial K$ is $C^2$, strictly convex, has positive curvature, and oriented counter-clockwise.
Parameterize it by the tangent angle $\alpha\in\mathbb R/2\pi\mathbb Z$:
\begin{equation}\label{eq:tangent-parametrization}
    \gamma'(\alpha)=\tau(\alpha)(\cos\alpha,\sin\alpha),
    \qquad ds=\tau(\alpha)\,d\alpha,
\end{equation}
where $\tau>0$ is the radius of curvature.  Put
$$
    P=\operatorname{per}K=\int_0^{2\pi}\tau(\alpha)\,d\alpha,
    \qquad
    S(t)=\int_t^{t+\pi}\tau(\alpha)\,d\alpha.
$$

For fixed $t$, choose a point uniformly with respect to arclength on the half-boundary whose tangent angles lie in $[t,t+\pi]$. If its tangent angle is $\alpha$, put $I_t = \alpha-t \in[0,\pi]$. Then $I_t$ has density
\begin{equation}\label{eq:one-sample-law}
    f_t(r)=\frac{\tau(t+r)}{S(t)}\,.
\end{equation}
Let $M_t$ be the median of three independent copies of $I_t$.

\begin{prop}\label{prop:random-separator}
With the preceding notation,
\begin{equation}\label{eq:median-reduction}
    \theta(K)
    \ge \frac{1}{6P}\int_0^{2\pi}
       S(t)^2\,\mathbb E\sin M_t\,dt.
\end{equation}
\end{prop}

\begin{proof}
Fix $t$ and project the body onto
$$
    u_t=(-\sin t,\cos t).
$$
The boundary $\partial K$ splits into two arcs joining the two support points with tangent vectors orthogonal to $u_t$. Let $\Gamma_t$ and $\Gamma_{t+\pi}$ be the two boundary arcs whose tangent angles range over $[t, t+\pi]$ and $[t+\pi, t+2\pi]$. Their arclengths are $S(t)$ and $S(t+\pi)$ = $P-S(t)$. Let $\mu_t, \mu_{t+\pi}$ be the pushforwards of normalized arclength measures on these arcs under projection onto $u_t$. Finally, let \(\nu_t\) be the pushforward of normalized arclength measure on the whole boundary \(\partial K\) under projection onto \(u_t\). Since
$$
\partial K=\Gamma_t\cup\Gamma_{t+\pi},
\qquad
S(t)+S(t+\pi)=P,
$$
we have
$$
\nu_t
=
\frac{S(t)}{P}\mu_t
+
\frac{S(t+\pi)}{P}\mu_{t+\pi}.
$$
As the Gini functional is concave under mixtures, 
\begin{equation}\label{concave}
\Gini(\nu_t)\ge\frac{S(t)}{P}\Gini(\mu_t)+\frac{S(t+\pi)}{P}\Gini(\mu_{t+\pi}).\end{equation}
Now
$$
|z|
=
\frac12\int_0^\pi
|\langle z,u_t\rangle|\,dt,
$$
and hence
$$
\theta(K)
=
\frac12\int_0^\pi \Gini(\nu_t)\,dt.
$$
Using \ref{concave},

$$\theta(K) \ge\frac1{2P}\int_0^\pi\left[S(t)\Gini(\mu_t)+S(t+\pi)\Gini(\mu_{t+\pi})\right]dt.$$
Projection of $\Gamma_{t+\pi}$ onto $u_{t+\pi}=-u_t$ reflects the
law obtained by projection onto $u_t$. Since the Gini mean difference
is invariant under reflection, the change of variable $r=t+\pi$ gives
\begin{equation}\label{theta_lower_bound}\theta(K) \geq \frac1{2P}\left[\int_0^\pi S(t)\Gini(\mu_t)\,dt+\int_\pi^{2\pi}S(t)\Gini(\mu_t)\,dt\right]=\frac1{2P}\int_0^{2\pi}S(t)\Gini(\mu_t)\,dt.
\end{equation}


Parametrize $\Gamma_t$ by arclength as
$$
r_t:[0,S(t)]\longrightarrow \Gamma_t,
$$
starting at the point with tangent angle $t$, and put
$$
x_t(s)=\langle r_t(s),u_t\rangle.
$$
If $\alpha_t(s)$ denotes the tangent angle at $r_t(s)$, then
$$
r_t'(s)=(\cos\alpha_t(s),\sin\alpha_t(s))
$$
and therefore
$$
x_t'(s)
=
\langle r_t'(s),u_t\rangle
=
\sin(\alpha_t(s)-t)\ge0.
$$
Thus $x_t$ is nondecreasing on $[0,S(t)]$.  Using that $x_t$ is
nondecreasing and applying Fubini's theorem we get
\begin{equation}\label{eq:chain-energy}
\begin{aligned}
    S(t)^2\Gini(\mu_t)
    &=\int_0^{S(t)}\int_0^{S(t)}
       |x_t(s)-x_t(r)|\,dr\,ds = 2\iint_{0\leq r\leq s \leq S(t)} x_t(s)-x_t(r)\,dr\,ds \\
    &=2\iint_{0\leq r\leq s \leq S(t)} \int_r^s x_t'(z)\,dz\,dr\,ds=2\int_0^{S(t)}z\bigl(S(t)-z\bigr)x_t'(z)\,dz.
\end{aligned}
\end{equation}
Observe that the median $Z_{(2)}$ of three independent uniform points $Z_1, Z_2,Z_3$ on $[0,S(t)]$ has density
$$
    \frac{6z(S(t)-z)}{S(t)^3}\,dz,
$$
so
$$S(t)^2\Gini(\mu_t) = \frac{S(t)^3}{3} \E x_t'(Z_{(2)}).$$ Since $s \mapsto \alpha_t(s)$ is increasing on the half-boundary, taking the median commutes with change of variable, i.e.
$$\alpha_t(Z_{(2)}) - t  = M_t$$

and, by \eqref{eq:tangent-parametrization}, $x_t'(s) = \sin(\alpha_t(s) - t)$.  Hence
$$    S(t)^2\Gini(\mu_t)=\frac{S(t)^3}{3}\,\mathbb E\sin M_t.$$

Finally, plugging this into \ref{theta_lower_bound} gives
\begin{equation}\label{tehta_lower_bound_fin}
\theta(K)
\ge
\frac1{6P}
\int_0^{2\pi}
S(t)^2\,\E\sin M_t\,dt.
\end{equation}
\end{proof}
Let us replace $M_t$ in the integral above with $I_t$. Then
$$
\begin{aligned}
\int_0^{2\pi}S(t)^2\,\E\sin I_t\,dt
&=
\int_0^{2\pi}
S(t)
\int_0^\pi
\tau(t+\omega)\sin \omega\,d\omega\,dt.
\end{aligned}
$$
The inner integral is the width $w(t)$ of $K$ in the direction $u_t$. Hence
\begin{equation}\label{second_brick}
\int_0^{2\pi}S(t)^2\E\sin I_tdt=\int_0^{2\pi}S(t)w(t)dt=\int_0^\pi\bigl(S(t)+S(t+\pi)\bigr)w(t)dt=P\int_0^\pi w(t)dt=P^2,
\end{equation}
where in the last step we used Cauchy's perimeter formula.
Therefore, the following Proposition together with Proposition~\ref{prop:random-separator} would conclude the proof of Theorem~\ref{thm:lower-boundary-distance} with non-strict inequality.
\begin{prop}\label{prop:median-concentration}
For every $0\le a\le\pi/2$,
\begin{equation}\label{eq:central-concentration}
\begin{aligned}
    \int_0^{2\pi}S(t)^2
       \mathbb P\{M_t\in[a,\pi-a]\}\,dt
    \ge
    \int_0^{2\pi}S(t)^2
       \mathbb P\{I_t\in[a,\pi-a]\}\,dt.
\end{aligned}
\end{equation}
Consequently,
\begin{equation}\label{eq:median-improves-sine}
    \int_0^{2\pi}S(t)^2\mathbb E\sin M_t\,dt
    \ge
    \int_0^{2\pi}S(t)^2\mathbb E\sin I_t\,dt.
\end{equation}
\end{prop}

For a randomly chosen half-boundary, the median of three random points tends to move away from its ends. If both angular end-pieces \([t,t+a]\) and \([t+\pi-a,t+\pi]\) contain at most half of the arclength, this already gives the statement by pointwise comparison in \(t\). The difficulty is that tangent angle and arclength need not be distributed uniformly: a short angular end-piece may carry more than half of the arclength. Proof of this Proposition relies on the following Lemma.
\begin{lem}\label{lem:cyclic-block}
Let $\mu$ be a finite nonnegative Borel measure on the circle of circumference $2h$.  Suppose
$0<a\le b$ and $a+b=h$.  For $x<y\le x+2h$, write $\mu[x,y)$ for the $\mu$-mass of the positively oriented half-open arc from $x$ to $y$. Put
\[
    \Phi(x,y)=\frac{xy(x-y)}{x+y}, \,\,\text{with} \,\,\Phi(0,0)=0,
\]
and define
\begin{equation}\label{eq:block-functional}
    \mathcal F_{a,b}(\mu)
    =\int_0^{2h}
      \Phi\bigl(\mu[t,t+a),\mu[t+a,t+a+b)\bigr)\,dt.
\end{equation}
Then
\begin{equation}\label{eq:block-inequality}
    \mathcal F_{a,b}(\mu)\le \mathcal F_{b,a}(\mu).
\end{equation}
\end{lem}
\begin{proof}
For $0<r<h$, put
\[
    L_r(t)=\mu[t,t+r),\qquad
    R_r(t)=\mu[t+r,t+h),\qquad
    H(t)=\mu[t,t+h).
\]
Thus
\[
    H(t)=L_r(t)+R_r(t),
\]
and
\[
    \Phi(L_r(t),R_r(t))
    =
    L_r(t)R_r(t)
    \left(
        1-2\frac{R_r(t)}{H(t)}
    \right).
\]
Whenever $H(t)=0$, we interpret $R_r(t)/H(t)$ as $0$; this is
harmless since then $L_r(t)R_r(t)=0$.

Since
\[
    L_r(t)R_r(t)
    =
    \int_{[t,t+r)}
    \int_{[t+r,t+h)}
    d\mu(y)\,d\mu(x),
\]
Fubini's theorem gives
\begin{equation}
\label{eq:block-pair}
    F_{r,h-r}(\mu)
    =
    \iint
    \left[
        \int_{\mathcal T_r(x,y)}
        \left(
            1-2\frac{R_r(t)}{H(t)}
        \right)dt
    \right]
    d\mu(x)\,d\mu(y),
\end{equation}
where
\[
    \mathcal T_r(x,y)
    =
    \left\{
        t:
        x\in[t,t+r),\
        y\in[t+r,t+h)
    \right\}.
\]

We shall call the decomposition

$$ [t,t+h)=[t,t+r)\cup[t+r,t+h) $$

the \(r\)-splitting of the half-arc. Thus \(F_{a,b}\) corresponds to the \(a\)-splitting and \(F_{b,a}\) to the \(b\)-splitting.
Fix an ordered pair $(x,y)$ and let $d$ be the positively oriented distance from $x$ to $y$. Choose the angular origin at \(x\), and represent \(y\) by the distance \(d\in[0,2h)\) from \(x\). Thus \(x=0\) and \(y=d\) on the periodic lift to \(\mathbb R\). 


The admissible values of \(t\), namely those for which \(0\in[t,t+r)\) and \(d\in[t+r,t+h)\), are therefore
\[
    T_r(d)
    =
    (-r,0]\cap(d-h,d-r].
\]
 Since $a+b=h$ and $a\le b$, one has
\[
\begin{aligned}
0\le d\le a:\qquad&
    T_a(d)=(-a,d-a],
    &
    T_b(d)=(-b,d-b],
\\
a\le d\le b:\qquad&
    T_a(d)=(-a,0],
    &
    T_b(d)=(d-h,d-b],
\\
b\le d\le h:\qquad&
    T_a(d)=(d-h,0],
    &
    T_b(d)=(d-h,0].
\end{aligned}
\]
Hence
\begin{equation}
\label{eq:block-translation}
    T_b(d)=T_a(d)-\delta(d),
\end{equation}
where
\[
    \delta(d)=
    \begin{cases}
        b-a, & 0\le d\le a,\\
        b-d, & a\le d\le b,\\
        0,   & b\le d\le h.
    \end{cases}
\]

If $h<d<2h$, then $T_r(d)=\varnothing$ for every $0<r<h$,
so there is nothing to prove. We may therefore assume $0\le d\le h$. Thus the admissible starting positions for the \(a\)- and \(b\)-splittings are paired by
\[
    t\longmapsto t'=t-\delta(d).
\]
In other words, for a fixed $t\in T_a(d)$ the cut point of $a$-splitting is $q=t+a$, whereas for the corresponding $b$-splitting it is $q' = t'+b = t+b-\delta(d)$. As $x$ and $y$ are fixed we write $\delta = \delta(d)$.
Note that $q\le q'$, since $\delta\le b-a$. Set
\[
    \lambda=\mu[t-\delta,t),\qquad
    C=\mu[t,q'),\qquad
    V=\mu[q',t+h-\delta),\qquad
    \eta=\mu[t+h-\delta,t+h).
\]
Then
\[
\begin{aligned}
    \frac{R_a(t)}{H(t)}
    &=
    \frac{\mu[q,t+h)}{\mu[t,t+h)}\ge
    \frac{\mu[q',t+h)}{\mu[t,t+h)}
    =
    \frac{V+\eta}{C+V+\eta}
\\
    &\ge
    \frac{V}{\lambda+C+V}
    =
    \frac{\mu[q',t+h-\delta)}
         {\mu[t-\delta,t+h-\delta)}=\frac{R_b(t')}{H(t')}.
\end{aligned}
\]
The second inequality follows by cross-multiplication:
\[
    (V+\eta)(\lambda+C+V)
    -
    V(C+V+\eta)
    =
    \lambda(V+\eta)+C\eta
    \ge0.
\]
Consequently,
\[
    1-2\frac{R_a(t)}{H(t)}
    \le
    1-2\frac{R_b(t')}{H(t')}.
\]
We defined $T_r$ so that $\mathcal T_r(x,y) = x+T_r(d)$. From \eqref{eq:block-translation} one gets \[
    \mathcal T_b(x,y)=\mathcal T_a(x,y)-\delta(d).
\]
Thus, for the fixed ordered pair $(x,y)$, the substitution $s=t-\delta(d)$ gives
\[
\int_{\mathcal T_b(x,y)}
\left(
    1-2\frac{R_b(s)}{H(s)}
\right)\,ds
=
\int_{\mathcal T_a(x,y)}
\left(
    1-2\frac{R_b(t-\delta(d))}
             {H(t-\delta(d))}
\right)\,dt
\ge
\int_{\mathcal T_a(x,y)}
\left(
    1-2\frac{R_a(t)}{H(t)}
\right)\,dt.
\]
% \[
%     \int_{\mathcal T_a(x,y)}
%     \left(
%         1-2\frac{R_a(t)}{H(t)}
%     \right)dt
%     \le
%     \int_{\mathcal T_b(x,y)}
%     \left(
%         1-2\frac{R_b(t)}{H(t)}
%     \right)dt.
% \]
Integrating this inequality over $(x,y)$ in
\eqref{eq:block-pair} gives
\[
    F_{a,b}(\mu)\le F_{b,a}(\mu).
\]
\end{proof}

Now we are ready to prove Proposition~\ref{prop:median-concentration}.
\begin{proof}[Proof of Proposition~\ref{prop:median-concentration}]
Denote tangent-angle measure of the body $K$ by $\mu$. That is, for a Borel set \(A\) of angles,
$\mu(A)=\mathcal H^1\left\{x\in \partial K: \alpha(x) \in A\right\}$, where $\alpha(x)$ is the tangent angle at $x$, as in \ref{eq:tangent-parametrization}, and $\mathcal H^1$ is the arclength measure. For fixed $t$, let
$$
    F_t(w)=\mathbb P(I_t<S w) = \frac{1}{S(t)}\mu[t,t+w),\quad S(t)=\mu[t,t+\pi).
$$
The CDF of $M_t$ is
$$
    G_t(w)=3F_t(w)^2-2F_t(w)^3,
$$
so
\begin{equation}\label{eq:median-cdf-difference}
    G_t(w)-F_t(w)
    =F_t(w)(1-F_t(w))(2F_t(w)-1).
\end{equation}
Put
$$
    A_t(w)=S(t)F_t(w)=\mu[t,t+w),
    \quad
    B_t(w)=S(t+w)F_{t+w}(\pi-w)=\mu[t+w, t+\pi),
$$
and then \eqref{eq:median-cdf-difference} says
\begin{equation}\label{eq:phi-cdf}
    S(t)^2\bigl(G_t(w)-F_t(w)\bigr)
    =\Phi\bigl(A_t(w),B_t(w)\bigr).
\end{equation}
After evaluating \eqref{eq:phi-cdf} at $w=a$ and $w=\pi-a$, the left-hand side of
\eqref{eq:central-concentration} minus the right-hand side equals
$$
    \mathcal F_{\pi-a,a}(\mu)-\mathcal F_{a,\pi-a}(\mu),
$$
which is nonnegative by Lemma~\ref{lem:cyclic-block}.  Finally,
\begin{equation}\label{eq:sine-layer-cake}
    \sin w=\int_0^{\pi/2}
       \mathbf 1_{[a,\pi-a]}(w)\cos a\,da,
    \qquad 0\le w\le\pi.
\end{equation}
Integrating \eqref{eq:central-concentration} against $\cos a\,da$ proves
\eqref{eq:median-improves-sine}.
\end{proof}





\begin{proof}[Proof of Theorem~\ref{thm:lower-boundary-distance}]
As we outlined before Proposition~\ref{prop:median-concentration},
\begin{equation}\label{eq:one-sample-baseline}
\begin{aligned}
    \int_0^{2\pi}S(t)^2\mathbb E\sin I_t\,dt
    &=\int_0^{2\pi}S(t)
      \left(\int_t^{t+\pi}\tau(\alpha)\sin(\alpha-t)\,d\alpha\right)dt\\
    &=\int_0^{2\pi}S(t)w(t)\,dt,
\end{aligned}
\end{equation}
where $w(t)$ is the width of $K$ in direction $u_t$. Indeed, $$w(t) = \langle \gamma(t+\pi) - \gamma(t), u_t \rangle = \int_t^{t+\pi}\langle \gamma'(\alpha), u_t \rangle d\alpha = \int_t^{t+\pi} \tau(\alpha)\sin(\alpha-t) d\alpha.$$ Since
$$
    S(t)+S(t+\pi)=P,
    \qquad
    w(t+\pi)=w(t),
$$
Cauchy's perimeter formula gives
\begin{equation}\label{eq:baseline-p-squared}
    \int_0^{2\pi}S(t)w(t)\,dt
    =P\int_0^\pi w(t)\,dt=P^2.
\end{equation}
Combining Propositions~\ref{prop:random-separator} and ~\ref{prop:median-concentration} with
\eqref{eq:baseline-p-squared}, we obtain
$$
    \theta(K)\ge \frac{P}{6}.
$$
By smooth approximation, the same chain of inequalities holds for every planar convex body, with \(\tau(\alpha)d\alpha\) replaced by its tangent-angle measure. Thus \(\theta(K)\ge P/6\) in general. It remains to exclude equality for a nondegenerate convex body.
Suppose therefore that
\[
    \theta(K)=\frac{P}{6}.
\]
Then both inequalities supplied by Propositions~5.2 and~5.3 must be
equalities.

Let $\mu$ be the tangent-angle measure of $\partial K$. We first
observe that equality in Proposition~5.3 implies that no open
semicircle contains three points of $\operatorname{supp} \mu$. Indeed, suppose that,
after lifting the angles to $\mathbb R$,
\[
    0=x<z<y=d<\pi,
    \qquad x,z,y\in\operatorname{supp}\mu.
\]
 Choose $a>0$ so small that
\[
    a<z,
    \qquad
    a<\pi-d,
\]
and put $b=\pi-a$.  Then $T_a(d) = (-a,0], \, \delta=b-d$. The cut points of $a$- and $b$-splittings are $q=t+a$ and $q'=t'+b=t+d$, respectively. Now, we want $z$ to lie strictly between them. For that purpose any $t$ in the non-empty interval $(\max\left\{-a, z-d\right\},0)$ works. For every such $t$ the inequality 
$$    \frac{\mu[q,t+h)}{\mu[t,t+h)}\ge
    \frac{\mu[q',t+h)}{\mu[t,t+h)}$$
is in fact strict, since $z \in [q,q')$. The same
remains true for $x$ and $y$ in sufficiently small neighborhoods of
the chosen points. Consequently,
\[
    F_{a,b}(\mu)<F_{b,a}(\mu).
\]
Moreover, this remains strict for all $a$ in a sufficiently small
open interval. That makes inequality \ref{eq:median-improves-sine} strict.
Now, we have shown that
\[
    |\operatorname{supp}\mu|\le4.
\]
Indeed, if five points occurred in cyclic order and $g_1,\ldots,g_5$
were the intervening angular gaps, then the absence of three points
in an open semicircle would imply
\[
    g_i+g_{i+1}\ge\pi
    \qquad (i\!\!\!\pmod 5).
\]
Summing gives $4\pi\ge5\pi$, since $\sum_i g_i=2\pi$, which is
impossible.

Hence $K$ is a polygon with at most four sides. We now show that
Proposition~5.2 is strict for every nondegenerate polygon. Choose an
open interval $J$ of angles $t$ so that for every $t\in J$
$$u_t = (-\sin t, \cos t)$$
is not normal to any edge of $K$.
For $t\in J$ the linear functional $x\mapsto\langle x, u_t \rangle$ attains its minimum and maximum at unique vertices $v$ and $w$, respectively. Those vertices split $\partial K$ in two complementary polygonal chains. Denote their lengths by $S_1,S_2$, and let $\nu_{1,t},\nu_{2,t}$ be their normalized arclength measures projected onto $u_t$.

The equality in the concavity step of
Proposition~5.2 is equivalent to
\[
    \nu_{1,t}=\nu_{2,t}.
\]

Let $e_1,e_2$ be the unit vectors along the edges of
the two chains incident to $v$ and directed away from $v$. For $t\in J$,
\[
    \langle e_i,u_t\rangle>0.
\]
Near the projection of $v$, the densities of
$\nu_{1,t}$ and $\nu_{2,t}$ are respectively
\[
    \frac{1}{S_1\langle e_1,u_t\rangle},
    \qquad
    \frac{1}{S_2\langle e_2,u_t\rangle}.
\]
Therefore equality of the two projected laws would imply
\[
    \bigl\langle S_1e_1-S_2e_2,u_t\bigr\rangle=0.
\]
As $S_1e_1-S_2e_2\ne0$, this can hold only for isolated values of $t$, whereas $J$ is an interval. Hence the Gini
concavity inequality is strict on a set of positive measure, and so
Proposition~5.2 is strict for every nondegenerate polygon.

Therefore,
\[
    \theta(K)>\frac{\operatorname{per}K}{6}.
\]
Sharpness follows from the rectangles
$K_\varepsilon=[0,1]\times[0,\varepsilon]$.
As $\varepsilon\downarrow0$, normalized boundary arclength concentrates
on the two horizontal sides, and hence
$\theta(K_\varepsilon)\to\mathbb E|U-V|=1/3$, where
$U,V\sim\operatorname{Unif}[0,1]$, while
$\operatorname{per}K_\varepsilon\to2$.
\end{proof}


\section{Note added}
 After the first version of our work was posted on arXiv, two related independent works became publicly available. Kukushkin \cite{Kukushkin_mach} gave a proof of Theorem~\ref{thm:endpoint} and the planar Zaporozhets--Tarasov conjecture using a quite different approach. Lotnikov \cite{Lotnikov_thes}, whose work is based on his bachelor's thesis defended in June 2026, constructed centrally symmetric counterexamples in every dimension $d\ge3$, for every moment of the distance. These results overlap, respectively, with our Theorem~\ref{thm:main} and the higher-dimensional counterexamples.

 \section{Acknowledgements}
 The author thanks A. Lotnikov and D. Zaporozhets for useful discussions.

 \section{Declaration of AI use}
 The author used ChatGPT with the GPT-5.6 Sol model during the preparation of
this work. The first proof of the planar case of Zaporozhets-Tarasov conjecture, namely Theorem~\ref{thm:main}, was obtained
by the author without the use of AI tools; it was based on KKT conditions and an
exhaustive case analysis and was substantially more technical. The argument was
subsequently simplified through discussions with ChatGPT. The author also used
ChatGPT to search for the explicit construction in Example 4.2 from the author's
underlying idea described in Remark~\ref{rmk:why} (that includes generation of relevant
code). The author reviewed and verified all AI-assisted output and takes full
responsibility for the content of the article.
\begin{thebibliography}{9}

\bibitem{BonnetEtAl}
G. Bonnet, A. Gusakova, C. Th\"ale, and D. Zaporozhets,
\emph{Sharp inequalities for the mean distance of random points in convex bodies},
Adv. Math. 386 (2021), 107813.

\bibitem{Lotnikov}
A. S. Lotnikov,
\emph{Mean distance between random points in centrally symmetric convex bodies},
Zap. Nauchn. Sem. POMI 544 (2025), 211--235.

\bibitem{TokmachevBoundary}
A. S. Tokmachev,
\emph{The Mean Distance Between Random Points on the Boundary of a Convex Figure},
J. Math. Sci. 286 (2024), 798--806.

\bibitem{Tokmachev2026}
A. S. Tokmachev,
\emph{Inequalities for convex functions of random points inside and on the boundary of convex bodies},
arXiv:2607.08869, 2026.

\bibitem{Khintchine}
A. Ya. Khintchine, \emph{On unimodal distributions},
Izv. Nauchno-Issled. Inst. Mat. Mekh. Tomsk. Gos. Univ. 2 (1938), 1–7. (In Russian.)
\bibitem{Yitzhaki}
S. Yitzhaki, \emph{Gini's mean difference: a superior measure of variability for non-normal distributions}, Metron 61 (2003), no. 2, 285–316.

\bibitem{HugSchneider} D. Hug and R. Schneider,
\emph{H\"older continuity for support measures of convex bodies},
Arch. Math. (Basel) 104 (2015), 83--92.
\bibitem{Schneider}
R. Schneider,\emph{ Convex Bodies: The Brunn–Minkowski Theory}, 2nd expanded ed., Cambridge University Press, Cambridge, 2014.

\bibitem{Kukushkin_mach}
M.~Kukushkin,
\newblock \emph{The Zaporozhets--Tarasov inequality for an arbitrary planar convex body},
\newblock arXiv:2608.14833, 2026.

\bibitem{Lotnikov_thes}
A.~S.~Lotnikov,
\newblock \emph{Mean distance between points inside and on the boundary of a convex body},
\newblock arXiv:2608.14848, 2026.
\end{thebibliography}

\end{document}
